% GENERATED FILE. Edit canonical lessons, then run npm run generate:books.
% canonical-source-hash: fnv1a64:cc84cd4264e39d01
% canonical-lessons: RU-C145-podpisyvat, RU-C145-pechatat, RU-C145-sokhranyat, RU-C145-skachivat, RU-C145-proigryvat

\chapter{To sign, To print, To keep, To download, To be beaten}
\label{ch:ru-to-sign-to-print-to-keep-to-download-to-be-beaten}
\hlchaptermodality{145}

\begin{chapteropening}
\textbf{By the end of this chapter:} \emph{I can say to sign, to print, to keep, to download and to be beaten.}

Complete the last lesson of chapter 145: I can say to sign, to print, to keep, to download and to be beaten.
\end{chapteropening}

\section[\texorpdfstring{\ru{подписывать} (podpísyvat)}{подписывать (podpísyvat)}]{\ru{подписывать} (podpísyvat) — to sign}

\label{lesson:RU-C145-podpisyvat}



\begin{quote}
Before the new one: say the Russian for to praise, then the Russian for to scold.
\end{quote}

\subsection*{You'll want to know: \ru{подписывать}}
\textbf{\ru{подписывать}} (\emph{podpísyvat}) — \textquotedblleft{}to sign\textquotedblright{}.

\begin{grammarlens}[title={What you've built}]
One more word for this chapter.
\end{grammarlens}

\subsection*{Your turn}
\begin{itemize}
\raggedright
  \item \emph{Say it:} \emph{podpísyvat}
  \item \emph{Say it:} \emph{podpísyvat}, once more
  \item \emph{Say it:} say \emph{podpísyvat}
  \item \emph{Recall:} say the Russian for to praise, then the Russian for to scold, then say \emph{podpísyvat} again
\end{itemize}

\begin{tcolorbox}[breakable,colback=teal!4,colframe=teal!35!black,title={Before you move on}]
What does \ru{подписывать} mean? (\textquotedblleft{}To sign\textquotedblright{}.) Say it once more.
\end{tcolorbox}
\section[\texorpdfstring{\ru{печатать} (pechátat)}{печатать (pechátat)}]{\ru{печатать} (pechátat) — to print, to type}

\label{lesson:RU-C145-pechatat}



\begin{quote}
Before the new one: say the Russian for to scold, then the Russian for to sign.
\end{quote}

\subsection*{You'll want to know: \ru{печатать}}
\textbf{\ru{печатать}} (\emph{pechátat}) — \textquotedblleft{}to print, to type\textquotedblright{}.

\begin{grammarlens}[title={What you've built}]
One more word for this chapter.
\end{grammarlens}

\subsection*{Your turn}
\begin{itemize}
\raggedright
  \item \emph{Say it:} \emph{pechátat}
  \item \emph{Say it:} \emph{pechátat}, once more
  \item \emph{Say it:} say \emph{pechátat}
  \item \emph{Recall:} say the Russian for to scold, then the Russian for to sign, then say \emph{pechátat} again
\end{itemize}

\begin{tcolorbox}[breakable,colback=teal!4,colframe=teal!35!black,title={Before you move on}]
What does \ru{печатать} mean? (\textquotedblleft{}To print, to type\textquotedblright{}.) Say it once more.
\end{tcolorbox}
\section[\texorpdfstring{\ru{сохранять} (sokhranyát)}{сохранять (sokhranyát)}]{\ru{сохранять} (sokhranyát) — to keep; to save (a file)}

\label{lesson:RU-C145-sokhranyat}



\begin{quote}
Before the new one: say the Russian for to sign, then the Russian for to print.
\end{quote}

\subsection*{You'll want to know: \ru{сохранять}}
\textbf{\ru{сохранять}} (\emph{sokhranyát}) — \textquotedblleft{}to keep; to save (a file)\textquotedblright{}.

\begin{grammarlens}[title={What you've built}]
One more word for this chapter.
\end{grammarlens}

\subsection*{Your turn}
\begin{itemize}
\raggedright
  \item \emph{Say it:} \emph{sokhranyát}
  \item \emph{Say it:} \emph{sokhranyát}, once more
  \item \emph{Say it:} say \emph{sokhranyát}
  \item \emph{Recall:} say the Russian for to sign, then the Russian for to print, then say \emph{sokhranyát} again
\end{itemize}

\begin{tcolorbox}[breakable,colback=teal!4,colframe=teal!35!black,title={Before you move on}]
What does \ru{сохранять} mean? (\textquotedblleft{}To keep; to save (a file)\textquotedblright{}.) Say it once more.
\end{tcolorbox}
\section[\texorpdfstring{\ru{скачивать} (skáchivat)}{скачивать (skáchivat)}]{\ru{скачивать} (skáchivat) — to download}

\label{lesson:RU-C145-skachivat}



\begin{quote}
Before the new one: say the Russian for to print, then the Russian for to keep; to save.
\end{quote}

\subsection*{You'll want to know: \ru{скачивать}}
\textbf{\ru{скачивать}} (\emph{skáchivat}) — \textquotedblleft{}to download\textquotedblright{}.

\begin{grammarlens}[title={What you've built}]
One more word for this chapter.
\end{grammarlens}

\subsection*{Your turn}
\begin{itemize}
\raggedright
  \item \emph{Say it:} \emph{skáchivat}
  \item \emph{Say it:} \emph{skáchivat}, once more
  \item \emph{Say it:} say \emph{skáchivat}
  \item \emph{Recall:} say the Russian for to print, then the Russian for to keep; to save, then say \emph{skáchivat} again
\end{itemize}

\begin{tcolorbox}[breakable,colback=teal!4,colframe=teal!35!black,title={Before you move on}]
What does \ru{скачивать} mean? (\textquotedblleft{}To download\textquotedblright{}.) Say it once more.
\end{tcolorbox}
\section[\texorpdfstring{\ru{проигрывать} (proígryvat)}{проигрывать (proígryvat)}]{\ru{проигрывать} (proígryvat) — to be beaten, to lose (a game)}

\label{lesson:RU-C145-proigryvat}



\begin{quote}
Before the new one: say the Russian for to keep; to save, then the Russian for to download.
\end{quote}

\subsection*{You'll want to know: \ru{проигрывать}}
\textbf{\ru{проигрывать}} (\emph{proígryvat}) — \textquotedblleft{}to be beaten, to lose (a game)\textquotedblright{}.

\begin{grammarlens}[title={What you've built}]
One more word for this chapter.
\end{grammarlens}

\subsection*{Your turn}
\begin{itemize}
\raggedright
  \item \emph{Say it:} \emph{proígryvat}
  \item \emph{Say it:} \emph{proígryvat}, once more
  \item \emph{Say it:} say \emph{proígryvat}
  \item \emph{Recall:} say the Russian for to keep; to save, then the Russian for to download, then say \emph{proígryvat} again
\end{itemize}

\begin{tcolorbox}[breakable,colback=teal!4,colframe=teal!35!black,title={Before you move on}]
What does \ru{проигрывать} mean? (\textquotedblleft{}To be beaten, to lose (a game)\textquotedblright{}.) Say it once more.
\end{tcolorbox}
